% Subsection 5.3.2: Few-shot CoT 边界分析

本节回答第四章矩阵设计的两个核心设问：\textbf{LoRA 与 CoT 是否正交可加}，\textbf{few-shot 的边际贡献是否依赖训练侧}。分析均基于 5.2 节两张主结果表的差分取数。

\subsubsection{训练 $\times$ 推理交互效应}
以 $(H8 - H7) - (H2 - H1)$ 与 $(H10 - H9) - (H3 - H1)$ 为核心两条 difference-in-differences，三种典型情形对应 4.1 节定义的三种判定形态：
\begin{itemize}
    \item \textbf{正交可加}（两条 DiD $\approx 0$）：LoRA 与 CoT 在 pass@$k$ 上独立叠加，工程上可分别优化；
    \item \textbf{LoRA 放大 CoT}（两条 DiD 显著 $> 0$）：LoRA 已把推理习惯内化进参数，与显式 CoT 推理形成正协同；
    \item \textbf{LoRA 吸收 CoT}（两条 DiD 显著 $< 0$）：LoRA 已覆盖了 CoT 提供的大部分增益，显式两阶段推理反而带来边际递减。
\end{itemize}
作为补充证据，$(H5 - H4) - (H2 - H1)$ 与 $(H6 - H4) - (H3 - H1)$ 进一步控制「训练侧是否监督 CoT 格式」这一变量：若这两项 DiD 与前两项在符号上一致，说明交互效应与训练侧是否显式监督 CoT 无关；若符号反转，则 CoT 监督成为触发正交/吸收的必要条件，需要单列讨论。

实测结果按 pass@1/5/10 分列如下：
\begin{itemize}
    \item $(H8-H7)-(H2-H1) = +5.0 / +4.0 / +3.4$\,pp（CoT 监督 LoRA × CoT-ZS 推理）；
    \item $(H10-H9)-(H3-H1) = +3.1 / +2.1 / +1.7$\,pp（CoT+FS 监督 LoRA × CoT-FS 推理）；
    \item $(H5-H4)-(H2-H1) = -0.1 / -0.4 / -0.4$\,pp（仅代码监督 LoRA × CoT-ZS 推理）；
    \item $(H6-H4)-(H3-H1) = -0.6 / -1.1 / -1.0$\,pp（仅代码监督 LoRA × CoT-FS 推理）。
\end{itemize}
前两条 DiD 显著大于 0 且 pass@1/5/10 同号、后两条 DiD 在所有 $k$ 上量级 $\le 1.1$\,pp，相对前两条 DiD 小约一个数量级，可视作 0 附近的弱负扰动——\textbf{交互效应不是「全局正交」也不是「全局放大」，而是受训练侧是否显式监督 CoT 调节}：当训练侧已监督 CoT 格式（H7--H10 族），LoRA 与 CoT 推理形成正协同；而仅做纯代码监督时（H4--H6 族），两者退化为近似正交可加，且在含 few-shot 的推理侧下伴随轻微负向扰动（$-1.0$\,pp 量级），暗示 CoT 推理示例与「仅代码监督」LoRA 之间存在轻度格式不匹配。工程侧结论：\textbf{若推理阶段确认走 CoT 两阶段，则训练阶段应同步注入 CoT 监督以触发协同放大}；若推理阶段固定 Direct，则 CoT 监督会带来格式开销（见 5.3.1 节 $\Delta_2$）。

\subsubsection{Few-shot 边际一致性}
Few-shot 的边际增益在不同训练侧上是否稳定，由以下两条对比给出直接证据：
\begin{itemize}
    \item \textbf{$\mathrm{H3} - \mathrm{H2}$（Base）vs $\mathrm{H6} - \mathrm{H5}$（LoRA\_code）}：两条「$k$ 例相对 $0$ 例」的 CoT 边际，训练侧均\textbf{未显式监督 CoT 格式}，用以验证 few-shot 是否为训练侧无关的结构性增益；
    \item \textbf{$\mathrm{H10} - \mathrm{H8}$（LoRA\_cot，$k$ 匹配）}：训练与推理的 $k$ 同步由 $0$ 提升到 $\texttt{FS\_K}$，用以衡量在监督已含 CoT 的前提下「训练与推理 few-shot 同步放大」的净收益。
\end{itemize}
若前两条边际在符号与量级上一致、第三条显著大于前两条，则结论为：few-shot 的基础增益来自推理侧示例本身，而在训练侧同步注入示例可进一步被模型吸收为参数知识；若前两条不一致，则说明 few-shot 在不同训练侧上起作用的机制不同，需结合具体配置进一步审视。

实测结果：
\begin{itemize}
    \item $H3-H2 = +4.2 / +3.7 / +3.3$\,pp（基座，pass@1/5/10）；
    \item $H6-H5 = +3.7 / +3.0 / +2.7$\,pp（\texttt{LoRA\_code}）；
    \item $H10-H8 = +3.6 / +2.9 / +2.5$\,pp（\texttt{LoRA\_cot} 训练-推理 $k$ 同步放大）。
\end{itemize}
前两条边际符号一致且量级差 $\le 0.6$\,pp（远小于 2\,pp 阈值），判定为\textbf{一致}——\textbf{few-shot 是训练侧无关的结构性增益，其推理侧贡献不依赖于训练侧是否监督 CoT 格式}。第三条边际（3.6\,pp）与前两条量级相当（差 $\le 0.6$\,pp），并未显著大于前两条，说明在监督已含 CoT 的前提下，「训练与推理 $k$ 同步放大」相对 zero-shot 监督的额外协同收益有限——这与 5.3.1 节 $\Delta_3=+1.3$\,pp（训练侧 few-shot 的弱内化）相互印证：few-shot 的主要作用仍发生在推理侧，训练侧贡献量级偏小。

\subsubsection{训练/推理 CoT 等价性}
为厘清「把 CoT 写进参数」与「在推理时显式走 CoT」两条路径是否互为替代，检查以下两组不等式：
\begin{itemize}
    \item $\mathrm{H9} \gtrsim \mathrm{H10}$：若成立，说明训练时充分内化 CoT 后，推理阶段不再需要显式两阶段也能得到相近甚至更好的结果；
    \item $\mathrm{H7} \gtrsim \mathrm{H8}$：与上一条在零示例场景下互为对照。
\end{itemize}
两组不等式同时成立则支持「训练侧等价替代」结论；若不等式方向相反，说明两条路径在不同子任务上各有优势，适合作为未来工作中的混合推理策略起点（见 6.2 节）。

实测：$H7=57.20\%$ vs $H8=63.90\%$，$H7-H8=-6.7$\,pp；$H9=58.50\%$ vs $H10=67.50\%$，$H9-H10=-9.0$\,pp。两组不等式方向均与 $H_i \gtrsim H_{i+1}$ 假设相反——\textbf{显式 CoT 两阶段推理始终优于训练侧内化 + Direct 推理}，差距在含 few-shot 的 $k$ 匹配档上进一步放大（−6.7\,pp $\to$ −9.0\,pp）。结论：\textbf{「CoT 进参数」与「CoT 进推理」并非互为替代，而是互补关系}——这一结论在 6.2.4 节实验扩展方向中将作为「混合推理策略」（在线根据预算与题型在 Direct / CoT 两条路径之间切换）的实证起点。

\subsubsection{$k$ 的选择灵敏度}
本文固定 \texttt{FS\_K=2}，对 $k$ 的灵敏度不在主实验范围内。但可从已有结果中间接推断灵敏度方向：若 H6 − H5 远大于 H3 − H2，则 $k$ 带来的增益在 LoRA\_code 上被进一步放大，提示未来工作尝试更大 $k$；若差距反向，则 $k$ 的进一步增大可能撞上 \texttt{MAX\_LENGTH} 或上下文衰减的边界，应优先做截断策略优化而非简单增大 $k$。